\clearpage\pagestyle{otMain}\setlength{\otBlocEspacement}{2.5mm}\fontsize{11}{15}\selectfont\corrigesChapitre{8}{Prendre une fraction d’une quantité}
\begin{blocLivre}{cor-8-1}\begin{corrige}{1}{Des parts simples}\demarche{Partager selon le dénominateur.}
$18:3=\res6$ et $32:4=\res8$.
\end{corrige}\end{blocLivre}
\begin{blocLivre}{cor-8-2}\begin{corrige}{2}{Deux tiers de douze}\demarche{Partager puis multiplier.}
$12:3=4$ puis $4\times2=\res8$ euros. $100:4=25$ puis $25\times3=\res{75}$ mètres.
\end{corrige}\end{blocLivre}
\begin{blocLivre}{cor-8-3}\begin{corrige}{3}{Deux chemins, un produit}\demarche{Choisir l’ordre qui donne des calculs ronds.}
$60:4=15$ puis $15\times3=\res{45}$. Ou $3\times60=180$ puis $180:4=\res{45}$ : les deux chemins mènent au même produit.
\end{corrige}\end{blocLivre}
\begin{blocLivre}{cor-8-4}\begin{corrige}{4}{Encadrer avant de calculer}\demarche{Comparer la fraction à 1 et à un demi.}
$\frac23<1$ : le résultat est plus petit que 12 ; $\frac23>\frac12$ : plus grand que 6. Calcul : $\res8$, bien entre 6 et 12.
\end{corrige}\end{blocLivre}
\begin{blocLivre}{cor-8-5}\begin{corrige}{5}{Des fractions d’heure}\demarche{Prendre la fraction de 60 minutes.}
$60:4=\res{15}$ min ; $15\times3=\res{45}$ min ; $60:6=\res{10}$ min.
\end{corrige}\end{blocLivre}
\begin{blocLivre}{cor-8-6}\begin{corrige}{6}{La jauge du réservoir}\demarche{Calculer la part puis le complément.}
$48:8=6$ puis $6\times3=\res{18}$ litres. Il manque $48-18=\res{30}$ litres, soit les $\frac58$ du réservoir.
\end{corrige}\end{blocLivre}
\begin{blocLivre}{cor-8-7}\begin{corrige}{7}{Retrouver le tout}\demarche{Revenir d’abord à une seule part.}
Le tout vaut $15\times3=\res{45}$ euros. Un cinquième fait $30:2=15$ km, le trajet $15\times5=\res{75}$ km.
\end{corrige}\end{blocLivre}
\begin{blocLivre}{cor-8-8}\begin{corrige}{8}{La classe de 28}\demarche{Partager en 7 puis compter 4 parts.}
$28:7=4$, puis $4\times4=\res{16}$ filles, et $28-16=\res{12}$ garçons. Contrôle : $16+12=28$.
\end{corrige}\end{blocLivre}
\begin{blocLivre}{cor-8-9}\begin{corrige}{9}{Plus grand que le tout ?}\demarche{Comparer la fraction à l’unité.}
$\frac55$ de 35 vaut $\res{35}$ : la quantité entière. $\frac{10}5$ de 35 vaut $35\times2=\res{70}$ : une fraction plus grande que 1 $\res{\text{agrandit}}$.
\end{corrige}\end{blocLivre}
\begin{blocLivre}{cor-8-10}\begin{corrige}{10}{L’étape de la fraction}\demarche{Repérer le tout de chaque fraction.}
Matin : $60:4=\res{15}$ euros ; reste 45. Après-midi : $\frac13$ de 45, soit $\res{15}$ euros ; reste $45-15=\res{30}$ euros. La seconde fraction agit sur le reste, pas sur 60.
\end{corrige}\end{blocLivre}
\begin{blocLivre}{cor-8-11}\begin{corrige}{11}{Le jardin potager}\demarche{Traiter chaque fraction du même tout.}
Légumes : $120:5=24$, puis $24\times2=\res{48}$ m². Fleurs : $120:10=12$, puis $12\times3=\res{36}$ m². Pelouse : $120-48-36=\res{36}$ m².
\end{corrige}\end{blocLivre}
\begin{blocLivre}{cor-8-12}\begin{corrige}{12}{Le prix soldé}\demarche{Calculer la part conservée puis l’écart.}
Nouveau prix : $60:4\times3=\res{45}$ euros. Économie : $60-45=\res{15}$ euros, soit le quart du prix.
\end{corrige}\end{blocLivre}
\begin{blocLivre}{cor-8-13}\begin{corrige}{13}{La collecte}\demarche{Passer par la part déjà remplie.}
Reçus : $600:6\times5=500$ livres. Il manque $600-500=\res{100}$ livres, soit $\frac16$ de l’objectif.
\end{corrige}\end{blocLivre}
\begin{blocLivre}{cor-8-14}\begin{corrige}{14}{Qui a raison ?}\demarche{Effectuer les deux chemins jusqu’au bout.}
Aya : $6\times2=\res{12}$. Noé : $36:3=\res{12}$. $\res{\text{Les deux ont raison}}$ : partager puis multiplier, ou multiplier puis partager, mènent au même résultat.
\end{corrige}\end{blocLivre}
\begin{blocLivre}{cor-8-15}\begin{corrige}{15}{Le tout mystère}\demarche{Travailler sur la fraction restante.}
Il lui reste $\frac14$ des économies : un quart vaut 12 euros, donc le tout valait $12\times4=\res{48}$ euros. Contrôle : $\frac34$ de 48 font 36, et $48-36=12$.
\end{corrige}\end{blocLivre}
